\chapter{Referências Normativas e Documentos de Apoio}

Este manual considera normas técnicas e documentos institucionais relacionados à apresentação de trabalhos acadêmicos. As normas devem ser consultadas diretamente quando houver dúvida específica de interpretação.

\section{Normas ABNT consideradas}

As designações a seguir foram conferidas diretamente nos exemplares normativos disponíveis ao projeto em 7 de agosto de 2026:

\begin{itemize}
  \item ABNT NBR 14724:2024, quarta edição, de 16 de dezembro de 2024, \emph{Informação e documentação --- Trabalhos acadêmicos --- Apresentação};
  \item ABNT NBR 6028:2021, segunda edição, de 18 de maio de 2021, \emph{Informação e documentação --- Resumo, resenha e recensão --- Apresentação};
  \item ABNT NBR 6027:2012, segunda edição, de 11 de dezembro de 2012, \emph{Informação e documentação --- Sumário --- Apresentação};
  \item ABNT NBR 6024:2012, \emph{Informação e documentação --- Numeração progressiva das seções de um documento --- Apresentação};
  \item ABNT NBR 10520:2023, segunda edição, de 19 de julho de 2023, \emph{Informação e documentação --- Citações em documentos --- Apresentação};
  \item ABNT NBR 6023:2025, terceira edição, de 21 de maio de 2025, \emph{Informação e documentação --- Referências --- Elaboração}.
\end{itemize}

O projeto também utiliza, conforme a modalidade e o elemento, as ABNT NBR 15287, 6022, 10719, 6034 e 15437, além das \emph{Normas de apresentação tabular} do IBGE. A edição, a designação oficial e a situação de vigência desses documentos devem ser confirmadas no catálogo autorizado antes da publicação institucional, pois os respectivos exemplares normativos não integram o acervo local auditado nesta revisão.

\section{Decisões institucionais}

Nos pontos em que a norma não define todos os detalhes de apresentação, o projeto validou decisões institucionais para preservar unidade visual e facilitar a vida acadêmica. Exemplos incluem formato de capa, apresentação da banca, modelo de data de aprovação, paginação, títulos, tratamento de listas e padrões de resumo.

Essas decisões estão registradas nos documentos de controle e incorporadas aos modelos do projeto. A publicação oficial do manual permanece sujeita ao fluxo institucional próprio, sem transformar novamente essas decisões em propostas.

\section{Documentos institucionais da UnDF}

O conteúdo institucional foi confrontado com os seguintes instrumentos:

\begin{itemize}
  \item Lei Complementar distrital nº 987, de 26 de julho de 2021;
  \item Decreto distrital nº 42.333, de 26 de julho de 2021;
  \item Estatuto da UnDF, aprovado pela Resolução nº 3/2022;
  \item Regimento Geral da UnDF, aprovado pela Portaria nº 4/2022;
  \item Resolução nº 12/2023, sobre elaboração e reestruturação de projetos pedagógicos de curso;
  \item Regimento da Graduação, aprovado pela Instrução nº 24, de 15 de abril de 2025;
  \item PPCs, regulamentos de curso e atos posteriores aplicáveis.
\end{itemize}

Os atos oficiais devem ser consultados em sua versão vigente no Sistema Integrado de Normas Jurídicas do Distrito Federal e nos canais oficiais da UnDF.

\section{Ética, animais e proteção de dados}

Também foram considerados como documentos de apoio:

\begin{itemize}
  \item Resolução CNS nº 466/2012, sobre pesquisas envolvendo seres humanos;
  \item Resolução CNS nº 510/2016, aplicável às pesquisas em Ciências Humanas e Sociais;
  \item Lei nº 11.794/2008 e normas do Concea, sobre uso científico de animais;
  \item Lei nº 13.709/2018, Lei Geral de Proteção de Dados Pessoais;
  \item Lei nº 9.610/1998, sobre direitos autorais;
  \item Guia orientativo da Autoridade Nacional de Proteção de Dados sobre tratamento de dados pessoais para fins acadêmicos e de pesquisa;
  \item orientações do Sistema CEP/Conep e das comissões institucionais competentes.
\end{itemize}

Os guias da Autoridade Nacional de Proteção de Dados possuem caráter orientativo e não substituem a legislação nem a avaliação institucional do caso concreto.

\section{Acessibilidade}

As recomendações de acessibilidade consideram o Modelo de Acessibilidade em Governo Eletrônico (eMAG), as Diretrizes de Acessibilidade para Conteúdo Web (WCAG) e boas práticas de documentos digitais. A UnDF ainda deve definir os requisitos mínimos aplicáveis aos TCCs e aos modelos institucionais.

\section{Atualização do manual}

Este manual deve ser atualizado sempre que houver:

\begin{itemize}
  \item alteração normativa relevante;
  \item nova decisão institucional;
  \item mudança em regulamento de curso;
  \item identificação de problema recorrente em bancas;
  \item necessidade de orientar novos tipos de trabalho;
  \item atualização de procedimentos de entrega ou arquivamento.
\end{itemize}

\caixaatencao{A versão preliminar deste manual deve passar por validação institucional antes de ser publicada como orientação institucional definitiva.}
